% Generated from validated Article Draft Result. Do not hand-edit; revise the structured draft instead.
\section{能力を測り、制御する層が追いつき始めた}
\noindent\textbf{強いモデルが増えるほど、「何が強いのかをどう比較するか」と「どの能力をどの条件で運用するか」が独立した問題になった。Chatbot Arena/MT-Benchと、OpenAI・Anthropicのalignment / governance面の取り組みを、evaluation layerとcontrol layerの形成として並べる。}\autocite{src-14f8acc117f1,src-1df7a91d31ec,src-25adcaa1c831,src-394ea6d17701,src-8f63dfa1697f,src-a1c8df89a3c9,src-b276b4af9d56,src-b2d10cbd609c,src-b6a5fdb87e9f}

Chatbot ArenaとMT-Benchは、モデル比較を個々の発表資料だけに閉じず、複数モデルを同じ評価surfaceへ載せる方向を押し出した。とくに人間の選好を含む比較は、単一の自動benchmarkとは異なる情報を与える一方、それ自体が絶対的な能力真値になるわけではない。\autocite{src-14f8acc117f1,src-1df7a91d31ec,src-b2d10cbd609c}


OpenAIとAnthropicの年内のalignment / governance関連の取り組みは、モデルの能力向上とは別に、利用・評価・scale-upの条件を定義するcontrol surfaceが必要だという問題設定を明示した。研究結果、内部評価、公開policyは同じ証拠種別ではないため、本稿では区別して扱う。\autocite{src-25adcaa1c831,src-394ea6d17701,src-8f63dfa1697f,src-a1c8df89a3c9,src-b276b4af9d56,src-b6a5fdb87e9f}


benchmark scoreやhuman preferenceは、task設計、prompt、比較対象、評価者構成などに依存する。Arena/MT-Benchを共通比較基盤として評価しても、そこから全用途の総合順位を導くことはしない。\autocite{src-14f8acc117f1,src-1df7a91d31ec,src-b2d10cbd609c}


年末に向かうにつれ、将来のより強い能力をどのthresholdで評価し、どの運用条件と結び付けるかが公開policyの論点として前面に出た。ここで重要なのは具体的な安全性が達成済みとみなすことではなく、capability growthとcontrol obligationを結び付ける枠組み自体が明示されたことである。\autocite{src-25adcaa1c831,src-394ea6d17701,src-8f63dfa1697f,src-a1c8df89a3c9,src-b276b4af9d56,src-b6a5fdb87e9f}


\begin{claimboundary}
評価結果と安全・alignment上の主張は、それぞれのproject/vendorが示した条件と範囲に帰属する。本稿はpolicyの存在や問題設定を一次資料で確認するが、有効性を独立に実証したとは扱わない。\autocite{src-25adcaa1c831,src-a1c8df89a3c9,src-b276b4af9d56,src-b6a5fdb87e9f}
\end{claimboundary}
